% ================================================================
%  第 1 章 权作序言 (原书 p.27-30, 边码 A1-A6)
%  源文: _work/src1/ch01.md
% ================================================================
\chapter{权作序言}

% ---------------- p.27 ----------------

\noindent\pmark 1986年1月30日\par

\origpage{A1} 只差一篇序言没写,就可以把《收获与播种》交给印刷厂了。我发誓,我是怀着世上最好的意愿,想写出点像样的东西来的。这一回,要写点合乎情理的东西。三四页,不能再多,但要写得有分量,用来引荐这头一千多页的庞然大物。要写点能“抓住”厌倦了的读者的东西,让他隐约看到:在这令人生畏的“一千多页”里,也许有些叫他感兴趣的东西(甚至,谁知道呢,与他有关的东西?)。“抓人”不是我的风格,不是。可这一回,我偏要破一次例!那位“疯到敢冒这个险的出版商”(敢出版这头显然无法出版的巨兽),总得勉强收回成本吧。

结果,没成。我尽力了。而且不止一个下午,不像我原先打算的那样三下两下就完事。到明天,整整三个星期了,我一直在写它,纸页越堆越高。写出来的东西,肯定不是人们体面地称之为“序言”的那种。又一次搞砸了,没治!我这把年纪,改不了了——我也不是那块料,不是叫卖的料,也不是让人叫卖的料。即便只为讨人喜欢(讨自己喜欢,也讨朋友喜欢……)。

写出来的,是一种长长的、带评注的“巡游”,穿过我作为数学家的作品。这巡游首先是给“门外汉”的——给那个“对数学一窍不通”的人。也是为我自己,我此前从未有过这样一趟巡游的闲情。一针引一线,我发现自己被引到一些向来留在未言说之中的东西面前,并且把它们说了出来。说来也巧,这些也正是我觉得最本质的东西——在我的工作中,在我的作品里。这些东西没有一点技术性。我这份天真的“传达”是否成功,由你评判——这尝试本身大概也有点疯。我的满足与快乐,将是让你感受到了它们。这些东西,我那些博学的同行里,许多人已经不再感受得到了。也许他们变得太博学、太有名望了。这往往让人失去与单纯而本质之物的接触。

在这趟“巡游一部作品”中,我也稍稍谈到我的一生。也在这里那里,稍稍谈到《收获与播种》讲的是什么。在“巡游”之后的那封“信”(注明去年五月)里,我还会更详细地再谈。这封信原本是写给我从前的学生们,以及数学界里我“昔日的朋友”的。但它同样没有一点技术性。任何读者,只要想通过一份“现场”的叙述,了解究竟哪些前因后果最终促使我写下《收获与播种》,都可以毫无困难地读它。比起“巡游”,它更能让你预先尝到“大世界”——数学的“大世界”——里某种气氛的味道。也(和“巡游”一样)尝到我的表达风格的味道,据说有点特别。还有借这风格表达出来的那种精神——这种精神同样不是人人都欣赏的。

\origpage{A2} 在“巡游”里,以及在《收获与播种》各处,我谈数学工作。这工作我熟悉,是第一手的。我所说的多数东西,对一切创造性工作、一切发现的工作,肯定都是真的。至少对所谓“智力”工作是真的——那种主要“靠脑子”和靠写来进行的工作。这样的工作,其标志是理解的绽放与成熟,即对我们正在探究之物的理解。% ---------------- p.28 ----------------
但举一个处于另一极的例子:爱的激情,同样也是发现的冲动。它为我们打开一种所谓“肉体的”认识,这种认识同样在更新、在舒展、在加深。这两种冲动——不妨说,一种驱动着工作中的数学家,另一种驱动着情人——彼此之近,远过于人们通常的猜度,也远过于人们愿意向自己承认的。我希望《收获与播种》的这些页面,能在你的工作中、在你每日的生活里,帮你感受到这一点。

在巡游中,主要谈的是数学工作本身。相反,对于这工作所处的境况,对于在真正的工作时间之外起作用的那些动机,我几乎一言不发。这有可能给人的形象——给我的形象,或给数学家、给一般“科学家”的形象——添上一层固然讨喜、却已走样的光晕。像是“伟大而高尚的激情”,没有任何修正。总之,合于那位大大的“科学神话”(S 请大写!)的路数。那英雄的、“普罗米修斯式”的神话,文人与学者们争先恐后地掉了进去(如今还在掉)。也许只有史家有时能抵挡这诱人的神话。实情是:在“科学家”的种种动机里——那些有时驱使他毫不吝惜地投入工作的动机——野心与虚荣所起的作用,与在任何别的行当里一样重大,而且几乎同样普遍。其形式或粗或细,因人而异。我绝不认为自己是个例外。读我这份见证,我希望,在这一点上不会留下任何疑问。

\origpage{A3} 也确实,最贪婪的野心也无力发现哪怕一条数学陈述,更无力证明它——正如它(譬如)无力“让人勃起”(就字面意义而言)。无论男女,使人“勃起”的绝不是野心,不是想要出风头、想要展示一种力量的欲望——那力量在此是性方面的——恰恰相反!而是对某种既强有力、又极真实、同时又极细微之物的敏锐知觉。人们可以称之为“美”,而美只是那东西的千张面孔之一。有野心,并不必然妨碍人偶尔感到一个存在、一件事物的美,这我同意。但可以肯定的是:让我们感到这美的,不是野心……

第一个发现并驾驭了火的人,正是和你我一样的人。绝不是人们以“英雄”“半神”之类名目设想的那种人。肯定,像你像我一样,他尝过焦虑的咬啮,也尝过那种虚荣的药膏——它让人忘掉咬啮。但在他“认识”火的那一刻,既没有恐惧,也没有虚荣。英雄神话里的真相就在于此。当神话被用来替我们遮掩事物的另一面——同样真实、同样本质的另一面——它就变得寡淡无味,变成了药膏。

我在《收获与播种》里的用意,是谈这两个方面——谈求知的冲动,也谈恐惧及其虚荣的解药。我相信自己“理解”、至少是认识这冲动及其本性。(也许有一天我会惊奇地发现,自己原来错得多么厉害……)但说到恐惧与虚荣,以及由此派生的对创造力的暗中阻塞,我很清楚自己并未深入这大谜的底。我也不知道,在我余下的岁月里,是否会看到这奥秘的底……

在写作《收获与播种》的过程中,浮现出两个形象,分别代表人类这场冒险的两个方面。那就是孩子(亦名“工匠”),以及老板。在即将开始的巡游里,要谈的几乎只是“孩子”。副题“孩子与母亲”里的,也是他。这名字,我希望,会在巡游途中渐渐明朗起来。

\origpage{A4} 而在反思其余的全部篇幅里,相反,占据前台的主要是老板。他可不是白当老板的! % ---------------- p.29 ----------------
更确切地说,不是一个老板,而是些相互竞争的企业的老板们。不过也确实:所有老板在本质上都彼此相像。而一旦开始谈老板,也就意味着要出现“坏蛋”了。在反思的第(I)部分(“愚妄与更新”,紧接本篇导论性的部分——即“四个乐章的前奏”——之后),那个“坏蛋”主要是我。在其后三个部分里,则主要是“别人”。轮着来!

这就是说,除了深刻的哲学思考与(毫无悔意的)“忏悔”,还会有“硫酸肖像”(借用我一位同行兼朋友的说法——他自己也挨了几下……)。更不必说那些大手笔、毫不含糊的“手术”。罗贝尔·若兰(Robert Jaulin)\fnote{1}{}{罗贝尔·若兰是老朋友了。我似乎明白,面对民族学界的建制,他处在一个“白狼”般的地位,与我在数学“上流社会”面前的处境略有相似。}向我保证(半开玩笑地):我在《收获与播种》里做的是“数学界的民族志”(也许是社会学,我也说不准了)。当然,听说自己(连自己都不知道)做着高深学问,总是受用的!事实是:在反思的“调查”部分里(而且违背我的意愿……),我看见我正在写下的那些纸页上,数学界建制的一大半人列队走过,更不必说许多地位更卑微的同行与朋友。而在最近这几个月里,自从去年十月我把《收获与播种》的试印本寄出去之后,这事又“来了一遍”。我这份见证,真是像一块石头扔进了死水塘。回响各色各样,几乎什么调子都有(只除了无聊那一种……)。几乎每一次,都完全出乎我的预料。此外还有许多沉默,意味深长。显然,关于我那些从前的学生、以及处境或高或低的同行们脑子里在转什么,我过去(如今也还)有许多要学,而且是形形色色的功课——请原谅,我是说“数学界的社会学”!对所有已经前来为我晚年这部社会学巨著添砖加瓦的人,我要在此表达我的感激之情。

当然,对那些暖色调的回响,我格外敏感。也有少数几位同行,向我道出了他们的感动,或一种(此前未表达过的)危机感、衰败感——对他们自认所属的这个数学界内部的危机与衰败。

\origpage{A5} 在这个圈子之外,最早对我的见证报以热烈、甚至激动回应的人里,我想在此举出西尔维·谢瓦莱与卡特琳·谢瓦莱(Sylvie et Catherine Chevalley)\fnote{2}{}{西尔维与卡特琳·谢瓦莱是克劳德·谢瓦莱的遗孀与女儿;这部《收获与播种》的核心部分(ReS III,“阴阳之钥”)正是献给他的。反思中多处谈到他,以及他在我的历程中所扮演的角色。}、罗贝尔·若兰、斯泰凡·德利若尔日(Stéphane Deligeorge)、克里斯蒂安·布尔茹瓦(Christian Bourgois)。如果说《收获与播种》将获得比最初那份(只给极狭小圈子的)试印本更广的流传,那主要归功于他们。尤其归功于他们那种有感染力的确信:我竭力去把握、去说出的东西,是应当被说出的;而且它能够被听到的范围,比我那些同行(往往阴沉,甚至乖戾,毫不愿意自我质疑……)的圈子更大。于是克里斯蒂安·布尔茹瓦毫不犹豫地承担了出版这“无法出版之物”的风险,而斯泰凡·德利若尔日则给我这份难以下咽的见证以殊荣,把它收进“Epistémé”丛书,(眼下)与牛顿、居维叶、阿拉戈并列。(我想不出更好的同伴了!)对每一位在我格外“敏感”的时刻一再表示同情与信任的人,我很高兴在此说出我全部的谢意。

于是我们这就出发,开始一趟巡游一部作品的旅程,作为穿越一生的旅行的开场。是的,一趟漫长的旅行,一千多页,每一页都塞得满满。这旅行我花了一辈子去做,也没做完,又花了一年多一页一页地重新发现它。有时词语迟迟不肯来,来表达一份仍在躲闪、尚未被把握的经验的全部汁液——就像压榨机里堆得满满的、熟透而饱满的葡萄,有时似乎想躲开紧攥着它的那股力量……但即便在词语仿佛互相推挤、汩汩涌出的时刻,它们的推挤与涌出也绝不是信手拈来。每一个词都在途中被称过,或者事后被称过,若发现太轻或太重,就仔细调准。所以这份反思—见证—旅行,不是为着让人一天读一点、一个月读完的,也不是为着那个急于读到“最后一句话”的读者。

% ---------------- p.30 ----------------
《收获与播种》里没有“最后一句话”,没有“结论”,正如我的一生里、你的一生里也没有。有的是一坛酒,在我这个人的木桶里陈了一辈子。你喝下的最后一杯不会比第一杯或第一百杯更好。它们全是“同一杯”,又各不相同。而倘若第一杯坏了,整桶都坏了;那还不如喝好水(若找得到),别喝坏酒。

\origpage{A6} 但好酒不宜匆匆喝,也不宜站着喝。
